\section{Introduction}

Ce chapitre situe le projet dans son cadre général. Il présente d'abord l'organisme d'accueil, puis le
domaine pharmaceutique dans lequel s'inscrit la solution. Il décrit ensuite le contexte du projet, les
solutions existantes, leurs limites et la réponse proposée à travers PharmacieConnect.

\section{Présentation de l'organisme d'accueil}

Le projet a été réalisé au sein de l'organisme d'accueil suivant : [À compléter].

\begin{table}[H]
    \centering
    \caption{Fiche signalétique de l'organisme d'accueil}
    \begin{tabular}{L{5cm}L{8cm}}
        \toprule
        \textbf{Élément} & \textbf{Information} \\
        \midrule
        Nom de l'organisme & [À compléter] \\
        Secteur d'activité & [À compléter] \\
        Adresse & [À compléter] \\
        Encadrant professionnel & [À compléter] \\
        Missions principales & [À compléter] \\
        \bottomrule
    \end{tabular}
\end{table}

Les informations détaillées sur l'organisme d'accueil doivent être complétées à partir de sources
vérifiées. Pour cette raison, le rapport ne contient pas de description non confirmée de l'entreprise, de
son organisation ou de ses services internes.

\section{Présentation du domaine pharmaceutique}

La pharmacie joue un rôle direct dans l'accès aux médicaments et dans l'orientation du patient. En dehors
des horaires habituels, les pharmacies de garde permettent de maintenir ce service disponible, notamment
la nuit, les week-ends et les jours fériés.

Dans ce domaine, la qualité des données a un impact réel sur l'utilisateur. Une adresse incorrecte, un
numéro de téléphone absent ou un horaire non actualisé peut empêcher une personne de trouver la
pharmacie recherchée. La mise à jour des informations devient donc une partie importante du service, au
même titre que leur consultation.

\section{Contexte du projet}

PharmacieConnect a été pensé comme un outil permettant de regrouper les informations relatives aux
pharmacies, aux plannings de garde et aux médicaments. La plateforme s'adresse à deux types
d'utilisateurs. Le citoyen utilise l'application mobile pour chercher une pharmacie, consulter les gardes
ou accéder au catalogue de médicaments. L'administrateur, de son côté, utilise le tableau de bord web
pour gérer et contrôler les données.

Le backend joue le rôle de point central entre ces deux interfaces et la base de données. Cette
organisation évite que les clients manipulent directement les données et permet de placer les règles
métier, la validation et la sécurité au niveau du serveur.

\section{Étude de l'existant}

Les informations sur les pharmacies de garde peuvent être consultées à travers plusieurs supports. Le
premier support est la liste statique, publiée sous forme de page web, d'image ou de document. Cette
solution est simple à diffuser, mais elle reste limitée lorsqu'il faut rechercher par ville, par date ou par
position.

Le deuxième support repose sur la recherche manuelle : appels téléphoniques, réseaux sociaux ou
informations locales. Cette méthode peut parfois donner une réponse rapide, mais elle dépend de la
disponibilité des personnes et de la fiabilité des informations partagées.

Le troisième support correspond aux applications ou sites spécialisés. Ces outils sont plus pratiques pour
l'utilisateur, surtout lorsqu'ils proposent la géolocalisation ou la recherche par date. En revanche, ils
nécessitent une base de données propre, une administration régulière et des mécanismes de mise à jour.

\begin{table}[H]
    \centering
    \caption{Comparaison générale des solutions existantes}
    \begin{tabular}{L{4cm}L{4.5cm}L{4.5cm}}
        \toprule
        \textbf{Solution} & \textbf{Avantages} & \textbf{Limites} \\
        \midrule
        Listes statiques & Faciles à publier, accessibles sans installation & Recherche difficile, mise à jour manuelle, faible adaptation mobile \\
        Appels ou sources locales & Possibilité d'obtenir une information directe & Dépendance à la disponibilité et risque d'information incomplète \\
        Applications spécialisées & Recherche plus rapide, filtres et géolocalisation possibles & Besoin d'une base fiable, d'une administration et d'une maintenance continue \\
        \bottomrule
    \end{tabular}
\end{table}

\section{Critique de l'existant}

L'étude de l'existant montre que le problème ne se limite pas à l'affichage des données. Les informations
peuvent être dispersées, présentées sous des formats différents ou mises à jour de manière irrégulière.
Pour un utilisateur mobile, cela rend la recherche moins fluide, surtout lorsqu'il souhaite filtrer par
gouvernorat, par date ou par proximité.

Le même problème se retrouve côté administration. Lorsque les données augmentent, la saisie manuelle
devient difficile à contrôler. Sans import, validation, gestion des rôles et journalisation, il devient plus
compliqué de garder une base cohérente et de suivre les modifications effectuées.

\section{Solution proposée}

La solution proposée, \textbf{PharmacieConnect}, regroupe trois composants :

\begin{itemize}
    \item une API backend développée avec FastAPI pour centraliser les données, gérer l'authentification,
    appliquer les règles métier et exposer les endpoints REST ;
    \item un tableau de bord administrateur développé avec React et Vite pour gérer les pharmacies, les
    gardes, les médicaments, les imports CSV, les assistants régionaux, les analytics et les journaux d'audit ;
    \item une application mobile développée avec Expo et React Native pour permettre au citoyen de
    rechercher des pharmacies, consulter les gardes, utiliser la géolocalisation et accéder au catalogue
    de médicaments.
\end{itemize}

Cette organisation répond aux limites relevées dans l'existant. Les données sont centralisées, les accès
sont contrôlés et les interfaces sont adaptées aux deux usages principaux : la consultation mobile et
l'administration.

\section{Conclusion}

Ce chapitre a présenté le contexte du projet, le domaine pharmaceutique, les solutions existantes et leurs
limites. PharmacieConnect propose une réponse structurée à ces limites en séparant le backend,
l'administration web et l'application mobile. Le chapitre suivant présente l'analyse détaillée des besoins
du système.
